\scanpage{7}
\mspage{6}
en fait, \deleted{et} décrivant $M(S)$. Pratiquement, on n'a guère qu'à considérer \deleted{dans} \deleted{\unsure{sur|une}} chaque question une ss-catégorie de type fini de $M(S)$, soit $C$. Elle est alors décrite par un \deleted{\unsure{groupe}} groupe algébrique \inserted{réductif} $G$ sur $Q$. Question intéressante : comment sont reliés ces groupes pour des $\xi$ différents (ils se déduisent tous les uns des autres par twist\ldots). Déjà intéressant si $S$ est le spectre d'une ext. finie $k$ de $Q$\ldots. !

Si $S$ est de car. $p>0$, on ne pourra en général trouver de foncteur fibre sur $Q$ (contre-exemple de Serre, lié aux courbes elliptiques de Hasse nul). Mais on aura un foncteur fibre sur $Q_\ell$ pour tt pt géom de car. $\ne\ell$, d'où un foncteur fibre sur une extension finie \inserted{$K$} de $Q$ (si on se restreint à une $C\subset M(S)$ stable par \deleted{\unsure{$\otimes$, fibre}} $\otimes,\vee$, et à \uline{\unsure{engendrt|engendrée}} fini\ldots). Donc on trouve un groupe algébrique réductif $G_K$ sur $K$, qui permet de représenter la catégorie des $K'$-motifs pour tt corps de base $K'$ qui est une extension de $K$.

\margin{////}
\inserted{Structure $\varepsilon,\unsure{\eta|D|y}$, \unsure{$\varepsilon\eta$\textup{ gr}$=\lambda^{?}$|$\varepsilon y(\lambda)=\lambda^2$}, structure $\unsure{\eta_1,\eta_2|\gamma_1,\gamma_2|y,i_1,i_2}$}
\deleted{\unsure{pour la compréhension du|pour corps de base} $G$, \unsure{$T^*$|$T_{\mathrm{Hdg}}$}}
\unsure{9)} \uline{Relations avec} \uline{le} $\pi_1$ \uline{\deleted{profini}\inserted{Hdg}}, : \inserted{A) $G/G^0\unsure{-|\longrightarrow}\pi^{1}$.} Soit $F$ un foncteur fibre à coeff$^{\mathrm{ts}}$ dans $K$ sur $M(C)$, $\xi$ un pt géométrique. On dit que $F$ et $\xi$ sont associés si on a un isom
% archive sheet number: 6

